% =======================================================
% ООП
% =======================================================

\section{ООП: абстракция, инкапсуляция, наследование, полиморфизм}

% -------------------------------------------------------
\subsection{Четыре столпа ООП}

\textbf{Объектно-ориентированное программирование} строится на четырёх принципах:

\begin{center}
\begin{tabular}{lp{9cm}}
\toprule
\textbf{Принцип} & \textbf{Суть} \\
\midrule
Абстракция     & выделить существенное, скрыть детали реализации \\
\addlinespace
Инкапсуляция   & объединить данные и методы, ограничить прямой доступ к данным \\
\addlinespace
Наследование   & строить новые типы на основе существующих \\
\addlinespace
Полиморфизм    & один интерфейс --- разное поведение для разных типов \\
\bottomrule
\end{tabular}
\end{center}

\subsection{Абстракция}

Абстракция --- предоставление \textbf{интерфейса} без раскрытия реализации.
Пользователь класса знает \textit{что} делает метод, но не \textit{как}:

\begin{lstlisting}
class Stack {
public:
    void push(int x);    // ЧТО делает -- интерфейс
    int  pop();
    bool empty() const;
private:
    // КАК устроено -- скрыто (вектор, список, массив -- неважно)
    std::vector<int> data_;
};
\end{lstlisting}

\subsection{Инкапсуляция}

Инкапсуляция --- сокрытие внутреннего состояния за модификаторами доступа. Доступ
к данным только через методы, которые поддерживают \textbf{инварианты} класса.

\begin{center}
\begin{tabular}{lp{8.5cm}}
\toprule
\textbf{Модификатор} & \textbf{Кто имеет доступ} \\
\midrule
\texttt{public}    & все \\
\texttt{protected} & сам класс и наследники \\
\texttt{private}   & только сам класс (и \texttt{friend}) \\
\bottomrule
\end{tabular}
\end{center}

\begin{lstlisting}
class Account {
public:
    void deposit(double sum) {
        if (sum > 0) balance_ += sum;   // инвариант: баланс корректен
    }
    double balance() const { return balance_; }
private:
    double balance_ = 0;   // нельзя изменить напрямую снаружи
};
\end{lstlisting}

\subsection{Наследование}

Наследование позволяет классу-наследнику (derived) получить члены базового
класса (base) и расширить/изменить их.

\begin{lstlisting}
class Animal {
public:
    void breathe() { /* общее поведение */ }
protected:
    std::string name_;
};
class Dog : public Animal {   // Dog -- это Animal (is-a)
public:
    void bark() { /* специфично для Dog */ }
};
\end{lstlisting}

\subsubsection{Виды наследования}

\texttt{public} наследование выражает отношение \textbf{«является» (is-a)} и
используется чаще всего. \texttt{protected}/\texttt{private} --- отношение
«реализован через», встречается редко.

\subsubsection{Множественное наследование и проблема ромба}

C++ допускает наследование от нескольких классов. Если два базовых класса сами
наследуют от общего предка, возникает \textbf{ромб (diamond)} --- дублирование
подобъекта. Решается \textbf{виртуальным наследованием}:

\begin{lstlisting}
class A { public: int x; };
class B : virtual public A {};   // virtual -> единственный подобъект A
class C : virtual public A {};
class D : public B, public C {}; // D::x -- один, а не два
\end{lstlisting}

\subsection{Полиморфизм}

\subsubsection{Динамический полиморфизм (виртуальные функции)}

Вызов метода через указатель/ссылку на базовый класс приводит к вызову версии
\textbf{реального} типа объекта --- в runtime:

\begin{lstlisting}
class Shape {
public:
    virtual double area() const = 0;   // чисто виртуальная (=0)
    virtual ~Shape() = default;        // виртуальный деструктор обязателен!
};
class Circle : public Shape {
    double r_;
public:
    Circle(double r) : r_(r) {}
    double area() const override { return 3.14159 * r_ * r_; }
};
class Square : public Shape {
    double s_;
public:
    Square(double s) : s_(s) {}
    double area() const override { return s_ * s_; }
};

void print_area(const Shape& s) {
    std::cout << s.area();   // вызовется нужная версия (Circle или Square)
}
\end{lstlisting}

\textbf{Абстрактный класс} --- содержит хотя бы одну чисто виртуальную функцию
(\texttt{= 0}); объект такого класса создать нельзя, только наследника.

\subsubsection{Как это работает: vtable и vptr}

У каждого полиморфного класса есть \textbf{таблица виртуальных функций (vtable)}
--- массив указателей на его методы. Каждый объект хранит скрытый указатель
\textbf{vptr} на vtable своего класса. Виртуальный вызов --- это косвенный вызов
через vtable:

\begin{lstlisting}[numbers=none]
Circle obj;            vtable Circle:
  +--------+           +------------------+
  | vptr   |---------> | Circle::area     |
  | r_     |           | Circle::~Circle  |
  +--------+           +------------------+
\end{lstlisting}

\textbf{Зачем виртуальный деструктор:} при \texttt{delete} через указатель на
базу без \texttt{virtual}-деструктора вызовется только \texttt{\textasciitilde Base},
ресурсы наследника утекут --- UB.

\subsubsection{\texttt{override} и \texttt{final}}

\begin{lstlisting}
class Derived : public Base {
    void f() override;   // компилятор проверит, что в базе есть virtual f
    void g() final;      // запретить переопределять g дальше
};
class X final {};        // запретить наследоваться от X
\end{lstlisting}

\subsubsection{Статический полиморфизм}

Полиморфизм без runtime-затрат достигается шаблонами и CRTP (см. соответствующие
разделы). Выбор реализации происходит на этапе компиляции.


% =======================================================
\section{Классы. Специальные методы}

% -------------------------------------------------------
\subsection{Специальные функции-члены}

Компилятор может сам сгенерировать шесть особых методов:

\begin{center}
\begin{tabular}{lp{8cm}}
\toprule
\textbf{Метод} & \textbf{Сигнатура} \\
\midrule
Конструктор по умолчанию & \texttt{T()} \\
Деструктор               & \texttt{\textasciitilde T()} \\
Копирующий конструктор   & \texttt{T(const T\&)} \\
Копирующее присваивание  & \texttt{T\& operator=(const T\&)} \\
Перемещающий конструктор & \texttt{T(T\&\&)} \\
Перемещающее присваивание & \texttt{T\& operator=(T\&\&)} \\
\bottomrule
\end{tabular}
\end{center}

\subsection{Конструкторы}

Конструктор инициализирует объект. Может быть перегружен.

\begin{lstlisting}
class Point {
    int x_, y_;
public:
    Point() : x_(0), y_(0) {}           // по умолчанию
    Point(int x, int y) : x_(x), y_(y) {}  // с параметрами
    explicit Point(int v) : x_(v), y_(v) {}  // explicit -- без неявных
};
\end{lstlisting}

\subsubsection{Список инициализации членов}

Члены инициализируются \textbf{в списке инициализации}, а не в теле конструктора.
Это важно: в теле было бы присваивание (поле уже создано), а в списке ---
прямая инициализация. Для \texttt{const}-полей и ссылок список обязателен.

\begin{lstlisting}
class Widget {
    const int id_;     // const можно инициализировать ТОЛЬКО в списке
    std::string name_;
public:
    Widget(int id, std::string name)
        : id_(id), name_(std::move(name)) {}  // порядок -- как объявлены поля!
};
\end{lstlisting}

\textbf{Важно:} члены инициализируются в порядке \textbf{объявления в классе}, а
не в порядке в списке инициализации.

\subsubsection{Делегирующие конструкторы}

\begin{lstlisting}
class Point {
    int x_, y_;
public:
    Point(int x, int y) : x_(x), y_(y) {}
    Point() : Point(0, 0) {}   // делегирует другому конструктору
};
\end{lstlisting}

\subsection{Деструктор}

Деструктор освобождает ресурсы при уничтожении объекта. Вызывается автоматически.
Основа RAII.

\begin{lstlisting}
class File {
    FILE* f_;
public:
    File(const char* name) : f_(fopen(name, "r")) {}
    ~File() { if (f_) fclose(f_); }   // ресурс освобождается сам
};
\end{lstlisting}

\subsection{Правило трёх / пяти / нуля}

\begin{itemize}[leftmargin=2em]
    \item \textbf{Правило трёх:} если нужен пользовательский деструктор,
          копирующий конструктор или копирующее присваивание --- скорее всего,
          нужны все три (класс владеет ресурсом).
    \item \textbf{Правило пяти:} добавьте move-конструктор и move-присваивание
          для эффективности.
    \item \textbf{Правило нуля:} лучше \textbf{вообще не писать} спецметоды ---
          используйте типы, которые сами управляют ресурсами
          (\texttt{std::string}, \texttt{std::vector}, умные указатели), тогда
          компилятор сгенерирует всё правильно.
\end{itemize}

\begin{lstlisting}
// Правило нуля: ничего не пишем -- всё генерируется корректно
class Person {
    std::string name_;
    std::vector<int> scores_;
};  // копирование, перемещение, удаление -- автоматически и правильно
\end{lstlisting}

\subsection{\texttt{= default} и \texttt{= delete}}

\begin{lstlisting}
class T {
public:
    T() = default;                       // явно просим сгенерировать
    T(const T&) = delete;                // запрещаем копирование
    T& operator=(const T&) = delete;
    T(T&&) = default;                    // разрешаем перемещение
};
\end{lstlisting}

\subsection{Статические члены}

\texttt{static}-член принадлежит \textbf{классу}, а не объекту --- общий для всех
экземпляров:

\begin{lstlisting}
class Counter {
    static int count_;     // объявление (одна копия на класс)
public:
    Counter()  { ++count_; }
    ~Counter() { --count_; }
    static int count() { return count_; }  // static-метод: без this
};
int Counter::count_ = 0;   // определение вне класса (или inline в C++17)
\end{lstlisting}

\subsection{\texttt{const}-методы и \texttt{mutable}}

\texttt{const}-метод обещает не менять объект; его можно вызывать на
\texttt{const}-объекте. \texttt{mutable}-поле можно менять даже в
\texttt{const}-методе (для кэшей, мьютексов):

\begin{lstlisting}
class Cache {
    mutable int cached_ = -1;   // можно менять в const-методе
public:
    int get() const {
        if (cached_ < 0) cached_ = compute();
        return cached_;
    }
};
\end{lstlisting}

\subsection{\texttt{friend}}

\texttt{friend}-функция или класс получает доступ к \texttt{private}-членам:

\begin{lstlisting}
class Vector2 {
    double x_, y_;
    friend Vector2 operator+(const Vector2&, const Vector2&);  // имеет доступ
};
\end{lstlisting}

\subsection{\texttt{this}}

\texttt{this} --- указатель на текущий объект внутри нестатического метода. Тип
--- \texttt{T*} (или \texttt{const T*} в const-методе).

\begin{lstlisting}
class Builder {
    int v_ = 0;
public:
    Builder& set(int v) { v_ = v; return *this; }  // fluent-интерфейс
};
Builder{}.set(1).set(2);   // цепочка вызовов
\end{lstlisting}


% =======================================================
\section{Перегрузка операторов}

% -------------------------------------------------------
\subsection{Зачем перегружать операторы}

Перегрузка операторов позволяет пользовательским типам вести себя как
встроенные: складываться через \texttt{+}, сравниваться через \texttt{==},
печататься через \texttt{<<}. Это синтаксический сахар: \texttt{a + b} ---
это вызов \texttt{operator+(a, b)}.

\subsection{Способы перегрузки: метод vs свободная функция}

\begin{lstlisting}
struct Complex {
    double re, im;

    // как метод: левый операнд -- *this
    Complex operator+(const Complex& o) const {
        return {re + o.re, im + o.im};
    }
};

// как свободная функция (симметрично для обоих операндов)
Complex operator-(const Complex& a, const Complex& b) {
    return {a.re - b.re, a.im - b.im};
}
\end{lstlisting}

\textbf{Правило:} бинарные симметричные операторы (\texttt{+}, \texttt{==}) лучше
делать свободными функциями --- тогда работают неявные преобразования для
\textit{обоих} операндов. Операторы, меняющие объект (\texttt{+=}, \texttt{=},
\texttt{[]}), --- методами.

\subsection{Арифметические операторы и каноничная форма}

Принято определять \texttt{+=} как метод, а \texttt{+} --- через него:

\begin{lstlisting}
struct Vec {
    double x, y;
    Vec& operator+=(const Vec& o) { x += o.x; y += o.y; return *this; }
};
Vec operator+(Vec a, const Vec& b) { return a += b; }  // через +=
\end{lstlisting}

\subsection{Операторы сравнения и spaceship (C++20)}

До C++20 нужно было определять все шесть операторов сравнения. В C++20 достаточно
оператора \texttt{<=>} (three-way comparison) --- остальные генерируются:

\begin{lstlisting}
#include <compare>
struct Point {
    int x, y;
    auto operator<=>(const Point&) const = default;  // <,>,<=,>=
    bool operator==(const Point&) const = default;   // ==, !=
};
\end{lstlisting}

\subsection{Оператор вывода \texttt{<<}}

Перегружается \textbf{только} как свободная функция (левый операнд ---
\texttt{std::ostream}, не наш тип):

\begin{lstlisting}
std::ostream& operator<<(std::ostream& os, const Point& p) {
    return os << '(' << p.x << ", " << p.y << ')';
}
std::cout << Point{1, 2};   // (1, 2)
\end{lstlisting}

\subsection{Оператор индексирования \texttt{[]}}

\begin{lstlisting}
class Array {
    int data_[100];
public:
    int&       operator[](size_t i)       { return data_[i]; }  // запись
    const int& operator[](size_t i) const { return data_[i]; }  // чтение
};
\end{lstlisting}

\subsection{Оператор вызова \texttt{()} --- функциональные объекты}

Класс с \texttt{operator()} называется \textbf{функтором} --- его объект можно
«вызывать»:

\begin{lstlisting}
struct Adder {
    int base;
    int operator()(int x) const { return x + base; }
};
Adder add5{5};
add5(10);   // 15 -- объект вызывается как функция
\end{lstlisting}

Функторы --- основа лямбд и предикатов для алгоритмов STL.

\subsection{Операторы инкремента}

Различают префиксный и постфиксный (отличаются фиктивным \texttt{int}-параметром):

\begin{lstlisting}
struct Counter {
    int v = 0;
    Counter& operator++()    { ++v; return *this; }       // префиксный ++c
    Counter  operator++(int) { Counter t = *this; ++v; return t; } // c++
};
\end{lstlisting}

\subsection{Операторы преобразования}

\begin{lstlisting}
struct Fraction {
    int num, den;
    explicit operator double() const { return double(num) / den; }
};
double d = static_cast<double>(Fraction{1, 2});  // 0.5
\end{lstlisting}

\subsection{Что нельзя перегружать}

Нельзя перегрузить: \texttt{::} (разрешение области), \texttt{.} (доступ к члену),
\texttt{.*}, \texttt{?:} (тернарный), \texttt{sizeof}, \texttt{typeid}. Нельзя
придумывать новые операторы или менять их арность/приоритет.

\subsection{Умные указатели и операторы \texttt{*}, \texttt{->}}

Перегрузка \texttt{operator*} и \texttt{operator->} позволяет классу вести себя
как указатель --- так устроены \texttt{unique\_ptr} и итераторы:

\begin{lstlisting}
template<typename T>
class Ptr {
    T* p_;
public:
    T& operator*()  const { return *p_; }
    T* operator->() const { return p_; }
};
\end{lstlisting}
